
% 依据：Expo12.pdf，局部页 1／合并页 757--758。
% 本组件延续至第 3 节。
\SGARefTarget{sga3:xii:expose}
\setcounter{footnote}{-1}
\footnotetext{2008 年 12 月 2 日 xy 版。}
从本报告起，与此前各报告不同，我们将使用代数闭域 \(k\) 上平滑代数群结构的已知结果，尤其是 1956/57 年 Chevalley 讨论班《代数 Lie 群的分类》〔\emph{Classification of Algebraic Lie Groups}〕所阐述的 Borel 仿射代数群理论。\footnote{编者注：此讨论班后来经 P.~Cartier 修订，作为 Chevalley 文集第 3 卷出版（Springer，2005）。}按照代数群理论中的惯例，我们用简称\emph{BIBLE} 引用这一讨论班。在以下各报告中，尤其需要\emph{BIBLE}~4、5、6 和 7 的结果（\emph{BIBLE} 后的数字为报告编号）。此外，看来概形理论并未给\emph{BIBLE} 中的 Borel 理论带来显著简化。因此，在本讨论班中重述该理论似无必要：这里的目的是从代数闭域上的已知结果，推导出任意基预概形上的类似结果。（Chevalley 的半单群结构理论则不同：看来在任意基预概形上\emph{从头}处理它更为有利。）
在口头报告中（第 1 至 4 节沿用其内容），我们限于 \(S\) 上的仿射群预概形，主要依靠报告 XI 第 4 节的可表性定理。在本笔记中（参见第 6 至 8 节），我们将简要说明如何用更简单的方法消去仿射性假设，而不使用报告 XI 中最精细的结果。本报告结果的进一步推广，另见报告 XV 和 XVI。（显然，报告 XI、XII、XV 和 XVI 的内容应当完全重写。）
\section{极大环面}
\SGARefTarget{sga3:xii:section:1}
\paragraph{\textbf{1.0.}}
\SGARefTarget{sga3:xii:1.0}
先设 \(G\) 为代数闭域 \(k\) 上的代数群。\emph{\(G\) 的极大环面}，是指 \(G\) 的代数子群 \(T\)，它是环面（这里 \(k\) 代数闭，即它同构于形如 \(\mathbf G_m^r\) 的群），且对此性质极大。注意，由于 \(k\) 完美，\(G_{\mathrm{red}}\) 是 \(G\) 的子群，在 \(k\) 上平滑，因此本质上是\emph{BIBLE} 意义下的代数群。此外，\(G\) 的每个约化子群（例如环面）自动是 \(G_{\mathrm{red}}\) 的子群。因此，\(G\) 的极大环面恰好是 \(G_{\mathrm{red}}\) 的极大环面。当 \(G\) 仿射，因而 \(G_{\mathrm{red}}\) 也仿射时，Borel 的一个基本定理断言，\(G\) 的任意两个极大环面都由
\[
  G(k)=G_{\mathrm{red}}(k)
\]
的某个元素共轭（\emph{BIBLE}~6，定理 4(c)）；特别地，它们维数相同。我们称 \(G\) 的极大环面的共同维数为 \emph{\(G\) 的约化秩}。实际上不必假设 \(G\) 仿射。这可由 Chevalley 的一个定理推出：\(k\) 上的每个连通平滑代数群，都是阿贝尔簇按仿射群的扩张；参见第 6 节。不过，在第 1 至 4 节中，我们多数时候仍限于在基上仿射的群预概形。
现在设 \(G\) 是 \(k\) 上的平滑代数群，\(T\) 是 \(G\) 的极大环面。称 \(T\) 在 \(G\) 中的中心化子 \(C(T)\) 为\emph{与 \(T\) 对应的 \(G\) 的 Cartan 子群}。对我们而言，它是 \SGARefLink{sga3:viii:corollary:6:7}{VIII~6.7} 所定义的群子预概形。不过，由于 \(G\) 在 \(k\) 上平滑，由 \SGARefLink{sga3:xi:corollary:5:3}{XI~5.3}，其 Cartan 子群 \(C\) 也在 \(k\) 上平滑。因此，此时 \(C\) 是 \(G\) 的唯一一个在 \(k\) 上平滑的群子预概形（即\emph{BIBLE} 意义下 \(G\) 的代数子群），使 \(C(k)\) 为 \(G(k)\) 中中心化 \(T(k)\) 的子群；换言之，它本质上就是\emph{BIBLE} 意义下的中心化子。
由刚才引用的共轭定理，不同极大环面对应的 Cartan 子群彼此共轭，因而维数相同。我们称其共同维数为 \emph{\(G\) 的幂零秩}；它等于 \(G_{\mathrm{red}}\) 的幂零秩。若 \(\rho_r(G)\) 与 \(\rho_n(G)\) 表示 \(G\) 的约化秩和幂零秩，则
\[
  \rho_r(G)\leq \rho_n(G),
\]
且当 \(G\) 仿射时，可将差
\[
  \rho_u(G)=\rho_n(G)-\rho_r(G)=\dim(C/T)
\]
称为 \emph{\(G\) 的幺幂秩}。若 \(G\) 平滑、仿射且连通，则 \(C\) 是连通幂零代数群（\emph{BIBLE}~6，定理 6(a)、(c)）。因此由\emph{BIBLE}~6，定理 2，它同构于积
\[
  C_s\times C_u,
\]
其中 \(C_s=T\) 是原来的极大环面（因而\emph{更是} \(C\) 的极大环面），而 \(C_u\) 是平滑幺幂子群，即由同构于 \(\mathbf G_a\) 的群逐次扩张得到的群（\emph{BIBLE}~6，定理 1、推论 1，以及\emph{BIBLE}~7，定理 4）。此时还有
\[
  \rho_u(G)=\dim C_u.
\]
\paragraph{\textbf{注 1.1。}}
\SGARefTarget{sga3:xii:1.1}
除刚才对仿射代数群给出的三种秩外，还有两种有用的秩。第一种是\emph{半单秩} \(\rho_s(G)\)，定义为商 \(G/R\) 的约化秩，其中 \(R\) 是 \(G\) 的根基。第二种是\emph{无穷小秩} \(\rho_i(G)\)，定义为 \(G\) 的 Lie 代数的幂零秩，下一报告将定义并研究它。本报告不使用它们，只记下不等式
\[
  \rho_s\leq \rho_r\leq \rho_n\leq \rho_i.
\]
\paragraph{\textbf{引理 1.2。}}
\SGARefTarget{sga3:xii:1.2}
\emph{设 \(G\) 是代数闭域 \(k\) 上的代数群，\(T\) 是 \(G\) 的代数子群，\(k'\) 是 \(k\) 的代数闭扩张。令 \(G'\) 和 \(T'\) 为 \(G\) 和 \(T\) 通过基扩张得到的群。则 \(T\) 是 \(G\) 的极大环面，当且仅当 \(T'\) 是 \(G'\) 的极大环面。}
这直接来自“有限扩张原理”（EGA~IV，9.1.1）。
\paragraph{\textbf{定义 1.3。}}
\SGARefTarget{sga3:xii:1.3}
设 \(S\) 是预概形，\(G\) 是有限型 \(S\)-群预概形，\(T\) 是 \(G\) 的群子预概形。若满足以下条件，则称 \(T\) 是 \emph{\(G\) 的极大环面}：
\begin{enumerate}
\renewcommand{\labelenumi}{\alph{enumi})}
\item \(T\) 是环面（\SGARefLink{sga3:ix:definition:1:3}{IX~1.3}）；且
\item 对每个 \(s\in S\)，若 \(\bar s\) 表示 \(\kappa(s)\) 的一个代数闭包的谱，则 \(T_{\bar s}\) 是 \(G_{\bar s}\) 的极大环面。
\end{enumerate}
\paragraph{\textbf{注 1.4。}}
\SGARefTarget{sga3:xii:1.4}
由 \SGARefLink{sga3:xii:1.2}{1.2}，当 \(S\) 是代数闭域的谱时，这恢复了通常的定义，且性质 \SGARefLink{sga3:xii:1.3}{1.3} 在任意基变换下稳定。注意，由 \SGARefLink{sga3:x:theorem:8-8}{X~8.8}，\SGARefLink{sga3:xii:1.3}{定义 1.3} 中的条件 (a) 可换为：
\begin{quote}
\textit{a\(^{\prime}\)) \(T\) 在 \(S\) 上有限呈示且平坦。}
\end{quote}
应当注意：\SGARefLink{sga3:xii:1.3}{定义 1.3} 意义下的极大环面，确实在 \(G\) 的子环面中极大，这由 \SGARefLink{sga3:ix:corollary:2:9}{IX~2.9} 立即得到；但反过来——例如基 \(S\) 连通时——只有 \(G\) 确实存在 \SGARefLink{sga3:xii:1.3}{1.3} 意义下的极大环面，逆命题才可能成立。一般并非如此，即便 \(S=\operatorname{Spec}(\mathbb Z)\) 且 \(G\)“半单”；另见 \SGARefLink{sga3:xii:1.6}{1.6}。不过，在报告 XIV 中将看到，当 \(S\) 是 Artin 概形，或 \(S\) 是局部概形而 \(G\)“约化”时，逆命题成立：此时 \(G\) 的每个环面都包含于一个极大环面。
\paragraph{\textbf{定义 1.5。}}
\SGARefTarget{sga3:xii:1.5}
设 \(G\) 是域 \(k\) 上的代数群。\(G\) 的\emph{约化秩}（分别为\emph{幂零秩}、\emph{幺幂秩}等等），是 \(G_{\bar k}\) 的约化秩（分别为相应的秩），其中 \(\bar k\) 是 \(k\) 的代数闭包。
由 \SGARefLink{sga3:xii:1.2}{1.2} 以及 \(\operatorname{Centr}_G(T)\) 的构造与基扩张相容，\SGARefLink{sga3:xii:1.5}{1.5} 引入的秩概念在基域扩张下不变。当 \(k\) 代数闭时，它们与本节开头引入的概念一致。
\paragraph{\textbf{注 1.6。}}
\SGARefTarget{sga3:xii:1.6}
不难在离散赋值环的谱 \(S\) 上构造平滑仿射群概形 \(G\)，使其一般纤维同构于 \(\mathbf G_m\)，特殊纤维同构于 \(\mathbf G_a\)。这样的 \(G\) 除退化为单位子群的平凡环面 \(T\) 外不含任何环面，而 \(T\) 显然不是极大环面。更确切地，在特殊纤维 \(G_0=\mathbf G_{a,k}\) 中，环面 \(T_0\) 确实极大；但在一般纤维 \(G_1=\mathbf G_{m,K}\) 中，环面 \(T_1\) 不再极大；这里 \(k\) 是剩余域，\(K\) 是赋值环的分式域。这个例子还说明，\(G_s\) 的约化秩（\(s\in S\)）不是 \(s\) 的连续函数。不过，有以下结果。
\paragraph{\textbf{定理 1.7。}}
\SGARefTarget{sga3:xii:1.7}
\emph{设 \(S\) 是预概形，\(G\) 是平滑仿射 \(S\)-群预概形。对每个 \(s\in S\)，置
\[
  \rho_r(s)=\rho_r(G_s),\qquad
  \rho_n(s)=\rho_n(G_s),
\]
即 \SGARefLink{sga3:xii:1.5}{1.5} 中 \(G_s\) 的约化秩和幂零秩。则：}
\begin{enumerate}
\renewcommand{\labelenumi}{\alph{enumi})}
\item\SGAInlineRefTarget{sga3:xii:1.7:item:a} \emph{\(S\) 上的函数 \(\rho_r\) 下半连续，而 \(\rho_n\) 上半连续。因此 \(\rho_u=\rho_n-\rho_r\) 上半连续。}
\item\SGAInlineRefTarget{sga3:xii:1.7:item:b} \emph{下列条件等价，且每个条件在任意基变换下稳定：}
\begin{enumerate}
\renewcommand{\labelenumii}{(\roman{enumii})}
\item \emph{\(S\) 上的函数 \(\rho_r\) 局部常值；}
\item \emph{\(G\) 对平展拓扑局部地存在极大环面；}
\item[\textup{(ii bis)}] \emph{\(G\) 对忠实平坦拟紧拓扑局部地存在极大环面。}
\end{enumerate}
\item \emph{设 \(T_1,T_2\) 是 \(G\) 的两个极大环面，因而特别满足 (b) 的条件。则 \(T_1\) 与 \(T_2\) 对平展拓扑局部共轭：存在满的平展态射 \(S'\to S\)，使 \(G_{S'}\) 的子群 \((T_1)_{S'}\) 和 \((T_2)_{S'}\) 由 \(G_{S'}\) 在 \(S'\) 上的一个截面共轭。}
\item \emph{在 (b) 的条件下，即 \(\rho_r\) 局部常值时，\(\rho_n\) 也局部常值，因而 \(\rho_u=\rho_n-\rho_r\) 亦然。}
\end{enumerate}
\emph{证明。}
a) 注意，对每个态射 \(S'\to S\)，若 \(G'=G\times_S S'\)，则由 \(G'\) 定义的 \(S'\) 上的函数 \(\rho'_r\) 等，与由 \(G\) 定义 \(\rho_r\) 等的方式相同，它们只须将后者与 \(S'\to S\) 复合即可得到。当 \(S'\to S\) 忠实平坦且拟紧时，\(\rho'\) 连续、上半连续或下半连续，当且仅当 \(\rho\) 分别具有相应性质，因为 \(S\) 的 Zariski 拓扑是 \(S'\) 的 Zariski 拓扑的商（SGA~1，\SGARefLink{sga3:viii:corollary:4:3}{VIII~4.3}）。因此，(a) 中的断言对忠实平坦拟紧拓扑是局部的。
设 \(s\in S\)。我们要证明，满足
\[
  \rho_r(t)\geq \rho_r(s)
  \qquad
  \bigl(\text{分别为 }\rho_n(t)\leq \rho_n(s)\bigr)
\]
的 \(t\in S\) 所组成的集合 \(U\) 是 \(s\) 的邻域。由有限扩张原理，存在 \(\kappa(s)\) 的有限扩张 \(k\)，使 \(G_k\) 存在极大环面。于是存在 \(s\) 的开邻域 \(U\) 及有限、满、局部自由态射 \(S'\to U\)，使纤维 \(S'_s\) 在 \(\kappa(s)\) 上同构于 \(\operatorname{Spec}(k)\)（参见 EGA~III，10.3.2，其中的 Noether 假设显然不必要）。由于 \(S'\to S\) 是开态射（SGA~1，\SGARefLink{sga3:iv:subsection:6:6}{IV~6.6}），可归约到 \(S'=S\)，即 \(G_s\) 存在极大环面 \(T_s\) 的情形。此外，由 \SGARefLink{sga3:xi:theorem:5:8:item:a}{XI~5.8(a)}，必要时再次用 \(S\) 上的平展 \(S'\) 替换 \(S\)，并取 \(s\) 上方的一点 \(s'\)，就可设 \(T_s\) 是 \(G\) 的某个子环面 \(T\) 在 \(s\) 处的纤维。于是对每个 \(t\in S\)，
\[
  \rho_r(t)=\rho_r(G_t)
  \geq \dim T_t
  =\dim T_s
  =\rho_r(G_s)
  =\rho_r(s),
\]
这就证明 \(\rho_r\) 下半连续。
另一方面，由 \SGARefLink{sga3:xi:corollary:5:3}{XI~5.3}，函子
\[
  \operatorname{Centr}_G(T)
\]
可由
\[
  C=\operatorname{Centr}_G(T),\footnote{编者注：此处修改了正文，以引入表示函子 \(\operatorname{Centr}_G(T)\) 的预概形 \(C=\operatorname{Centr}_G(T)\)。}
\]
表示，这是 \(G\) 的闭群子预概形，在 \(S\) 上平滑。因此存在 \(s\) 的邻域 \(U\)，使得 \(t\in U\) 蕴含
\[
  \dim C_t=\dim C_s=\rho_n(s).
\]
现在，\(\rho_n\) 的上半连续性来自关系
\[
  \rho_n(t)\leq \dim C_t
  \qquad\text{对每个 }t\in S,
\]
而这一关系包含在以下纯几何引理中。
\paragraph{\textbf{引理 1.8。}}
\SGARefTarget{sga3:xii:1.8}
\emph{设 \(k\) 是域，\(G\) 是 \(k\) 上的平滑仿射代数群，\(T\) 是 \(G\) 中的环面，\(C\) 是其中心化子。则}
\[
  \rho_n(G)\leq \dim C.
\]
事实上，可设 \(k\) 代数闭，于是 \(T\) 包含于极大环面 \(T'\)。若 \(C'\) 是 \(T'\) 的中心化子，则 \(C'\subset C\)，故 \(\dim C'\leq \dim C\)。引理得证。
b) 若 \(\rho_r\) 局部常值，则对 \(G\) 中的每个环面 \(T\) 及每个 \(s\in S\)，若 \(T_s\) 是 \(G_s\) 的极大环面，则存在 \(s\) 的开邻域 \(U\)，使 \(T|_U\) 是 \(G|_U\) 的极大环面。应用 (a) 的论证，可见 (i) 蕴含 (ii bis)。反之，(ii bis) 蕴含 (i)，因为当 \(G\) 存在极大环面 \(T\) 时，显然
\[
  \rho_r(s)=\dim T_s
\]
是 \(s\) 的局部常值函数；而在 (a) 的证明开头已指出，\(\rho_r\) 的连续性对忠实平坦拓扑是局部的。因此 (i) 等价于 (ii bis)，而 (ii) 显然蕴含 (ii bis)；只须再证明逆向蕴含 (i) \(\Rightarrow\) (ii)。
为此，引入 \SGARefLink{sga3:xi:theorem:4:1}{XI~4.1} 的函子 \(F\)，\footnote{编者注：这是 \(G\) 的乘法型子群函子。}它是在 \(S\) 上平滑且分离的预概形，并考虑 \(F\) 的子函子 \(\mathcal T\)，它在 \(S\) 上的 \(S'\) 处取值为 \(G_{S'}\) 中极大环面的集合。根据前面的观察，在 (i) 下，\(G_{S'}\) 中的环面若在一点 \(s\in S'\) 的纤维中极大，则在 \(s\) 的某个开邻域上极大，所以函子 \(\mathcal T\) 可由 \(F\) 的一个开子预概形表示，因而在 \(S\) 上平滑且分离。由于结构态射 \(\mathcal T\to S\) 显然满，由 \SGARefLink{sga3:xi:proposition:1:10}{XI~1.10}，它对平展拓扑局部地有截面。这证明了 (i) \(\Rightarrow\) (ii)。
c) 结合本节开头回顾的 Borel 共轭定理，这直接来自 \SGARefLink{sga3:xi:corollary:5:4:bis}{XI~5.4 bis}。
d) 根据 (a) 证明开头的观察以及 (b)，可设 \(G\) 存在极大环面 \(T\)。若 \(C\) 是其中心化子，则由 \SGARefLink{sga3:xi:corollary:5:3}{XI~5.3}，\(C\) 可表且在 \(S\) 上平滑，所以函数
\[
  s\longmapsto \rho_n(s)=\dim C_s
\]
确实局部常值。
\SGARefLink{sga3:xii:1.7}{1.7} 的证明至此完成。在这种情形，我们将说 \(G\) \emph{具有局部常值约化秩}，以指 \SGARefLink{sga3:xii:1.7:item:b}{1.7(b)} 的条件。记下如下推论。
\paragraph{\textbf{推论 1.9。}}
\SGARefTarget{sga3:xii:1.9}
\emph{设 \(G\) 如 \SGARefLink{sga3:xii:1.7}{1.7}，\(s\in S\) 满足 \(\rho_u(s)=0\)，即 \(\rho_r(s)=\rho_n(s)\)（等价地，\(G_{\bar k}\) 的 Cartan 子群都是环面，其中 \(\bar k\) 是 \(\kappa(s)\) 的代数闭包）。则存在 \(s\) 的开邻域 \(U\)，使 \(\rho_r\) 和 \(\rho_n\) 在其上常值。特别地，对每个 \(t\in U\)，\(G_t\) 的幺幂秩 \(\rho_u(t)\) 为零。}
事实上，这立即来自 \SGARefLink{sga3:xii:1.7:item:a}{1.7(a)} 以及对每个 \(t\in S\) 成立的不等式 \(\rho_r(t)\leq \rho_n(t)\)。
还应注意，在 (b) 的证明中，我们证明了以下结果。
\paragraph{\textbf{推论 1.10。}}
\SGARefTarget{sga3:xii:1.10}
\emph{设 \(G\) 如 \SGARefLink{sga3:xii:1.7}{1.7}，并设 \(G\) 具有局部常值约化秩。考虑函子
\[
  \mathcal T:(\mathbf{Sch})^\circ\longrightarrow(\mathbf{Ens})
\]
使得对 \(S\) 上的每个 \(S'\)，
\[
  \mathcal T(S')=
  \{G_{S'}\text{ 的极大环面}\}.
\]
则 \(\mathcal T\) 可由 \(S\) 上平滑、分离且有限型的预概形表示。}
只须验证 \(\mathcal T\) 在 \(S\) 上有限型。由于问题对忠实平坦拟紧拓扑是局部的，可设 \(G\) 存在极大环面 \(T\)。由 \SGARefLink{sga3:xi:corollary:5:3:bis}{XI~5.3 bis} 以及 \SGARefLink{sga3:xi:remark:5:5}{XI~5.5} 的 bis 版本，\(\operatorname{Norm}_G(T)\) 和 \(G/\operatorname{Norm}_G(T)\) 可分别由预概形 \(\operatorname{Norm}_G(T)\) 和 \(G/\operatorname{Norm}_G(T)\) 表示，而 \(\mathcal T\) 同构于
\[
  G/\operatorname{Norm}_G(T).\footnote{编者注：此处修改了正文，以引入预概形 \(\operatorname{Norm}_G(T)\) 和 \(G/\operatorname{Norm}_G(T)\)。}
\]
由 \(g\mapsto\operatorname{int}(g)(T)\) 定义的态射 \(G\to\mathcal T\) 是满的。由于 \(G\) 在 \(S\) 上拟紧，\(\mathcal T\) 也拟紧，证明完毕。
\paragraph{\textbf{注 1.11。}}
\SGARefTarget{sga3:xii:1.11}
a) \SGARefLink{sga3:xii:1.10}{1.10} 的预概形 \(\mathcal T\) 将恰当地称为 \emph{\(G\) 的极大环面预概形}。第 5 节将说明它实际上在 \(S\) 上仿射。当 \(G\) 同构于形如 \(\operatorname{GL}(n)\) 的预概形的闭群子预概形时，可直接验证这一点：使用 \SGARefLink{sga3:xi:remark:4:6}{XI~4.6}，并注意，由 \SGARefLink{sga3:xii:1.7:item:b}{1.7(b)} 证明中的论证，\(\mathcal T\) 可与表示 \(G\) 的乘法型子群的预概形 \(F\) 中一个既开又闭的子预概形等同。
b) 不必使用报告 XI 的精细结果，仅用初等结果 \SGARefLink{sga3:xi:proposition:3:12}{XI~3.12} 和 \SGARefLink{sga3:xi:theorem:6:2}{XI~6.2}，且只处理在基上有限的乘法型群（实质上是群 \({}_nT\)，其中 \(T\) 是极大环面），就可以证明 \SGARefLink{sga3:xii:1.7}{1.7}，从而证明 \SGARefLink{sga3:xii:1.9}{1.9}；比较第 7 节。同一观察也适用于 \SGARefLink{sga3:xii:1.10}{1.10} 的证明。
本节最后给出若干极大环面唯一的例子。
\paragraph{\textbf{命题 1.12。}}
\SGARefTarget{sga3:xii:1.12}
\emph{设 \(G\) 是有限型乘法型 \(S\)-群预概形。则 \(G\) 存在唯一的极大环面，且 \(G\) 中的每个环面都包含于该极大环面。}
唯一性显然来自最后一个断言，它将极大环面刻画为 \(G\) 的最大子环面。唯一性也使存在性成为对忠实平坦拟紧拓扑局部的问题。因此，可设 \(G\) 可对角化，例如 \(G=D_S(M)\)，其中 \(M\) 是有限生成交换群。令 \(M_0\) 为 \(M\) 对其挠子群的商。我们断言，\(G\) 中的环面
\[
  T=D_S(M_0)
\]
是极大环面，也是最大子环面。事实上，\(G\) 的子环面 \(T'\) 对 fpqc 拓扑局部可对角化。因此，为证明 \(T'\subset T\)，可设 \(T'\) 可对角化，因而形如 \(D_S(N)\)，其中 \(N\) 是 \(M\) 的自由商（\SGARefLink{sga3:viii:corollary:1:4}{VIII~1.4} 和 \SGARefLink{sga3:viii:corollary:3:2:item:b:01}{3.2(b)}）。可知 \(N\) 是 \(M_0\) 的商，故 \(T'\subset T\)。由于 \(T=D_S(M_0)\) 的构造与任意基扩张相容，这同时证明 \(T\) 是 \(G\) 的极大环面，证明完毕。当 \(G\) 在 \(S\) 上平滑时，\SGARefLink{sga3:xii:1.12}{命题 1.12} 有如下推广。
\paragraph{\textbf{命题 1.13。}}
\SGARefTarget{sga3:xii:1.13}
\emph{设 \(G\) 是在 \(S\) 上有限呈示的 \(S\)-群预概形。假设对忠实平坦拟紧拓扑局部地，\(G\) 存在中心极大环面。则 \(G\) 在整体上存在唯一的极大环面，\footnote{它是中心的！}且它是 \(G\) 的最大子环面。}
这里，唯一性仍然直接来自最后一个断言，并使存在性成为对 fpqc 拓扑局部的问题。因此可设 \(G\) 存在中心极大环面 \(T\)。我们要证明 \(G\) 的每个环面 \(R\) 都包含于 \(T\)。这来自下一个引理。
\paragraph{\textbf{引理 1.14。}}
\SGARefTarget{sga3:xii:1.14}
\emph{设 \(G\) 是在 \(S\) 上有限呈示的 \(S\)-群预概形，\(T\) 是 \(G\) 的极大环面，\(R\) 是 \(G\) 的子环面，并设 \(T\) 与 \(R\) 交换。则 \(R\subset T\)。}
事实上，由于 \(R\) 与 \(T\) 交换，态射
\[
  R\times_S T\longrightarrow G,\qquad (r,t)\longmapsto r\cdot t,
\]
是群同态。因此，由 \SGARefLink{sga3:ix:theorem:6:8}{IX~6.8}，它在 \(G\) 中有一个像子群 \(T'\)，后者为 \(R\times_S T\) 的乘法型商群，因而是环面，且显然包含 \(T\)。由于 \(T\) 是极大环面，由 \SGARefLink{sga3:xii:1.4}{1.4} 有 \(T'=T\)，故 \(R\subset T\)。这证明了 \SGARefLink{sga3:xii:1.14}{1.14}，从而证明了 \SGARefLink{sga3:xii:1.13}{1.13}。特别地，使用 \SGARefLink{sga3:xii:1.7:item:b}{1.7(b)} 得到：
\paragraph{\textbf{推论 1.15。}}
\SGARefTarget{sga3:xii:1.15}
\emph{设 \(G\) 是具有局部常值约化秩的交换、平滑、仿射 \(S\)-群预概形。则 \(G\) 存在唯一的极大环面，且此环面包含 \(G\) 的每个子环面。}
\paragraph{\textbf{推论 1.16。}}
\SGARefTarget{sga3:xii:1.16}
\emph{设 \(G\) 是平滑仿射 \(S\)-群预概形。假设对每个 \(s\in S\)，若 \(\bar s\) 表示 \(\kappa(s)\) 的一个代数闭包 \(\bar k\) 的谱，则几何纤维 \(G_{\bar s}\) 是连通幂零代数群（按\emph{BIBLE} 的意义，即其 \(\bar k\)-值点群 \(G(\bar k)\) 幂零）。再设 \(G\) 具有局部常值约化秩。则 \(G\) 存在唯一极大环面 \(T\)；此外，\(T\) 是中心的，且为 \(G\) 的最大子环面。}
由 1.7(b)，\(G\) 对 fpqc 拓扑局部地存在极大环面，而由 \SGARefLink{sga3:xii:1.13}{1.13}，只须证明 \(G\) 的极大环面是中心的。由 \SGARefLink{sga3:ix:theorem:5:6:item:b}{IX~5.6(b)}，可归约到 \(S\) 是域的谱的情形，并可设该域代数闭。于是，由\emph{BIBLE}~6，\SGARefLink{sga3:xii:section:2}{定理 2}，\(T(k)\) 位于 \(G(k)\) 的中心。这蕴含 \(T\) 位于 \(G\) 的中心，因为 \(\operatorname{Centr}_G(T)\) 是 \(G\) 的闭子概形（VIII~6.6），包含 \(G(k)\) 的所有点，而 \(G\) 约化，所以由零点定理，它等于 \(G\)。
最后，为便于以后引用，记下如下初等命题，前文已隐含地用到它。
\paragraph{\textbf{命题 1.17。}}
\SGARefTarget{sga3:xii:1.17}
\emph{设 \(G\supset H\supset T\) 是 \(S\) 上有限型群预概形。下列条件等价：}
\begin{enumerate}
\renewcommand{\labelenumi}{(\roman{enumi})}
\item \emph{\(T\) 是 \(G\) 的极大环面；}
\item \emph{\(T\) 是 \(H\) 的极大环面，且对每个 \(s\in S\)，约化秩满足等式
\[
    \rho_r(G_s)=\rho_r(H_s).
  \]
}
\end{enumerate}
\section{Weyl 群}
\SGARefTarget{sga3:xii:section:2}
先设 \(k\) 是代数闭域，\(G\) 是 \(k\) 上的平滑仿射代数群。若 \(T\) 是极大环面，\(C\) 是其中心化子，\(N\) 是其正规化子，则由 \SGARefLink{sga3:xi:proposition:5:9}{XI~5.9}，它们是 \(G\) 的平滑闭子群，且 \(C\) 是 \(N\) 的开子群。因此
\[
  W=N/C
\]
是 \(k\) 上的有限平展群，因而由其 \(k\)-值点群 \(W(k)\) 确定：
\[
  W=W(k)_k
\]
（即通常的有限群 \(W(k)\) 所定义的常值群）。称有限群 \(W(k)\) 为 \(G\) 相对于 \(T\) 的\emph{几何 Weyl 群}，在不致混淆时简称 \emph{Weyl 群}。由 Borel 共轭定理，不同极大环面对应的 Weyl 群彼此同构；因此，有时不指定极大环面，而称为 \(G\) 的“那个”Weyl 群。由于 \(C\)、\(N\) 和 \(N/C\) 的构造与任意基扩张交换，若 \(k'\) 是 \(k\) 的代数闭扩张，则 \(G_{k'}\) 相对于 \(T_{k'}\) 的几何 Weyl 群，典范同构于 \(G\) 相对于 \(T\) 的几何 Weyl 群。因此，\(G\) 的“那个”几何 Weyl 群——严格地说是通常有限群的一个同构类——与 \(G_{k'}\) 的几何 Weyl 群一致。
于是，当 \(G\) 是任意域 \(k\) 上的平滑仿射代数群时，可以将 \(G_{k'}\) 的几何 Weyl 群的同构类称为 \emph{\(G\) 的几何 Weyl 群}，其中 \(k'\) 是 \(k\) 的任意代数闭扩张。当 \(G\) 存在极大环面 \(T\) 时，可以如上构造
\[
  C=\operatorname{Centr}_G(T),\qquad
  N=\operatorname{Norm}_G(T),\qquad
  W=N/C
\]
后者是 \(k\) 上的有限平展群，称为 \emph{\(G\) 相对于 \(T\) 的 Weyl 群}。几何 Weyl 群就是 \(W\) 在 \(k\) 的任意代数闭扩张 \(k'\) 中取值的点群的类。一般而言，仅知道几何 Weyl 群 \(W(k')\)，显然不再足以重建代数群 \(W\)：还必须知道 \(k\) 在 \(k'\) 中的可分代数闭包的 Galois 群在 \(W(k')\) 上的作用。
最后，设 \(G\) 是任意基 \(S\) 上的群预概形，在 \(S\) 上平滑且仿射，\(T\) 是 \(G\) 的极大环面。由 \SGARefLink{sga3:xi:proposition:5:9}{XI~5.9}，仍可构造群
\[
  W(T)=N(T)/C(T),
\]
它是平展、分离、拟有限的 \(S\)-群预概形。其相对于剩余域 \(\kappa(s)\)（\(s\in S\)）的代数闭包的几何纤维，是纤维 \(G_s\) 的几何 Weyl 群。因此，\SGARefLink{sga3:xi:corollary:5:10}{XI~5.10} 给出了这些群随 \(s\in S\) 变化的信息。可将这些信息进一步精确化并推广如下。
\paragraph{\textbf{定理 2.1。}}
\SGARefTarget{sga3:xii:2.1}
\emph{设 \(S\) 是预概形，\(G\) 是平滑仿射 \(S\)-群预概形。对每个 \(s\in S\)，令 \(w(s)\) 为 \(G_s\) 的几何 Weyl 群，因而是有限群的同构类。在这些类的集合 \(E\) 上引入如下预序：\(w\leq w'\) 当且仅当 \(w\) 与 \(w'\) 可分别由有限群 \(W\) 与 \(W'\) 表示，使 \(W\) 同构于 \(W'\) 的某个子群的商。在此约定下：}
\begin{enumerate}
\renewcommand{\labelenumi}{\alph{enumi})}
\item\SGAInlineRefTarget{sga3:xii:2.1:item:a} \emph{从 \(S\) 到 \(E\) 的函数 \(s\mapsto w(s)\) 下半连续。}
\item\SGAInlineRefTarget{sga3:xii:2.1:item:b} \emph{设 \(G\) 具有局部常值约化秩。则以下条件等价：}
\begin{enumerate}
\renewcommand{\labelenumii}{(\roman{enumii})}
\item \emph{函数 \(s\mapsto w(s)\) 局部常值；}
\item[\textup{(i bis)}] \emph{函数 \(s\mapsto\operatorname{card}w(s)\) 局部常值；}
\item \emph{对平展拓扑局部地（或仅对 fpqc 拓扑局部地），存在极大环面 \(T\)，使 \(W(T)\) 在其基上有限；}
\item[\textup{(ii bis)}] \emph{对 \(S\) 上的每个 \(S'\) 以及 \(G_{S'}\) 的每个极大环面 \(T\)，对应的 Weyl 群 \(W(T)\) 在 \(S'\) 上有限。}
\end{enumerate}
\end{enumerate}
\emph{证明。}
a) 仿照 \SGARefLink{sga3:xii:1.7:item:a}{1.7(a)} 的论证，为证明对每个 \(s\in S\)，存在 \(s\) 的开邻域 \(U\)，使 \(t\in U\) 蕴含 \(w(t)\geq w(s)\)，可归约到 \(G\) 中存在环面 \(R\)，且 \(R_s\) 是 \(G_s\) 的极大环面的情形。置
\[
  W(R)=\operatorname{Norm}_G(R)/\operatorname{Centr}_G(R).
\]
与 \SGARefLink{sga3:xi:proposition:5:9}{XI~5.9} 一样，这是 \(S\) 上平展、分离、拟有限的群预概形。对 \(t\in S\)，令 \(w'(t)\in E\) 为它在 \(t\) 处的几何纤维。由于 \(R_s\) 是 \(G_s\) 的极大环面，且 \(\operatorname{Norm}\)、\(\operatorname{Centr}\) 和 \(\operatorname{Norm}/\operatorname{Centr}\) 的构造与任意基扩张交换，特别与取纤维交换，有
\[
  w(s)=w'(s).
\]
此外，我们断言，对靠近 \(s\) 的 \(t\)，
\[
  w(t)\geq w'(t)
  \qquad\text{且}\qquad
  w'(t)\geq w'(s),
\]
这就足以证明 (a)。这两个不等式包含在以下两个引理中。
\paragraph{\textbf{引理 2.2。}}
\SGARefTarget{sga3:xii:2.2}
\emph{设 \(S\) 是预概形，\(W\) 是平展、分离、拟有限的 \(S\)-群预概形。对每个 \(s\in S\)，令 \(f(s)\) 为 \(W\) 在 \(s\) 处的几何纤维的类，它是 \SGARefLink{sga3:xii:2.1}{2.1} 所引入的有限群同构类的有序集 \(E\) 的元素。则函数 \(f:S\to E\) 下半连续。它在 \(s\in S\) 的某个邻域内常值，当且仅当函数 \(s\mapsto\operatorname{card}f(s)\) 在那里常值；又当且仅当 \(W\) 在 \(s\) 的某个开邻域 \(U\) 上方于 \(S\) 上有限。}
这个结果以群论语言精细化了 \SGARefLink{sga3:xi:corollary:5:10}{XI~5.10} 证明中所用的结果。其证明完全标准，这里仅给出梗概。照例可归约到 \(S\) 仿射且 Noether 的情形。立即可见，函数 \(f\) 可构造（EGA \(0_{\mathrm{III}}\)，9.3.1、9.3.2，以及 EGA~IV，9 的一系列结果），而 EGA \(0_{\mathrm{III}}\)，9.3.4 将半连续性归约为证明：若 \(t\) 是 \(s\) 的一般化，则
\[
  f(t)\geq f(s).
\]
由 EGA~II，7.1.7，这归约到 \(S\) 是离散赋值环的谱的情形，并可设该环完备且剩余域代数闭。令 \(s\) 为 \(S\) 的闭点。由于 \(W\) 在 \(S\) 上平展且分离，它包含既开又闭且在 \(S\) 上有限的子概形 \(W'\)，使 \(W'_s=W_s\)（EGA~II，6.2.6）；这里立即可见 \(W'\) 是 \(W\) 的子群。此外，由于 \(W'\) 在 \(S=\operatorname{Spec}(V)\) 上有限平展，\(V\) 完备且剩余域代数闭，它是常值群，故形如 \(A_S\)，其中
\[
  A=W'(\kappa(s))=W(\kappa(s))
\]
的类为 \(f(s)\)。若 \(B\) 是 \(W\) 在 \(S\) 的一般点 \(t\) 处的几何纤维，则有典范单态射 \(A\to B\)，这证明了 \(f(s)\leq f(t)\)。注意，这个论证实际上证明了相对于 \(E\) 上比 \SGARefLink{sga3:xii:2.1}{2.1} 所指定的序更细的序的半连续性。
\(f\) 在 \(s\) 处常值与 \(s\mapsto\operatorname{card}f(s)\) 在该处常值的等价性，来自以下事实：对 \(w,w'\in E\)，
\[
  w\leq w'
  \quad\text{且}\quad
  \operatorname{card}w=\operatorname{card}w'
  \quad\Longrightarrow\quad
  w=w'.
\]
此条件与 \(W\) 在 \(s\) 的某个邻域上有限的等价性，与 \(W\) 的群结构无关；\SGARefLink{sga3:xi:corollary:5:10}{XI~5.10} 后已指出这一点，也容易由上述论证结合固有性的赋值判别法（EGA~II，7.3.8）得到。
\paragraph{\textbf{引理 2.3。}}
\SGARefTarget{sga3:xii:2.3}
\emph{设 \(G\) 是代数闭域 \(k\) 上的平滑仿射代数群，\(R\subset T\) 是两个子环面。令 \(W(R)\) 和 \(W(T)\) 为 \SGARefLink{sga3:xi:proposition:5:9}{XI~5.9} 所给的两个对应有限群，均为正规化子对中心化子的商。则 \(W(R)\) 同构于 \(W(T)\) 的某个子群的商。}
事实上，考虑下图
\[
\begin{tikzcd}[column sep=large,row sep=large]
T \arrow[r, hook] &
C(T) \arrow[r, hook] \arrow[d, hook] &
C(R) \arrow[d, hook] \\
& N(T)\cap N(R) \arrow[r, hook] \arrow[d, hook] &
N(R) \\
& N(T) &
\end{tikzcd}
\]
则 \((N(T)\cap N(R))/C(T)\) 是
\[
  W(T)=N(T)/C(T),
\]
的子群，并有显然的同态
\[
  (N(T)\cap N(R))/C(T)
  \longrightarrow
  W(R)=N(R)/C(R).
\]
一切归结为证明最后这个同态满，即对 \(N(R)\) 的每个 \(k\)-值点 \(g\)，存在 \(C(R)\) 的 \(k\)-值点 \(c\)，使 \(cg\) 正规化 \(T\)，或等价地，使
\[
  \operatorname{int}(c)\bigl(\operatorname{int}(g)T\bigr)=T.
\]
为此，只须注意，\(\operatorname{int}(g)T\) 是 \(N(R)\) 的环面，因而是 \(C(R)\) 的环面，因为后者是 \(N(R)\) 的开子群。环面 \(T\) 和 \(\operatorname{int}(g)T\) 是 \(C(R)\) 的极大环面，因为它们在 \(G\) 中极大，故断言来自 Borel 共轭定理。
这证明了 \SGARefLink{sga3:xii:2.3}{2.3}，从而证明了 \SGARefLink{sga3:xii:2.1:item:a}{2.1(a)}。
b) 已经指出，(i) 与 (i bis) 显然等价。由 \SGARefLink{sga3:xii:2.2}{2.2} 的逆向部分，或等价地由 \SGARefLink{sga3:xi:corollary:5:10}{XI~5.10}，它们蕴含 (ii bis)。由 \SGARefLink{sga3:xii:1.7:item:b}{1.7(b)}，条件 (ii bis) 蕴含 (ii)。最后，(ii) 蕴含 (i bis)：与 \SGARefLink{sga3:xii:1.7:item:a}{1.7(a)} 一样，条件 (i) 对 fpqc 拓扑是局部的，所以可设 \(G\) 存在极大环面 \(T\)，使 \(W(T)\) 在 \(S\) 上有限，再次由 \SGARefLink{sga3:xi:corollary:5:10}{XI~5.10} 即得结论。\SGARefLink{sga3:xii:2.1}{2.1} 证毕。
\section{Cartan 子群}
\SGARefTarget{sga3:xii:section:3}
\paragraph{\textbf{定义 3.1。}}
\SGARefTarget{sga3:xii:3.1}
设 \(G\) 是预概形 \(S\) 上的平滑有限型群预概形。\emph{\(G\) 的 Cartan 子群}是指 \(G\) 的群子预概形 \(C\)，它在 \(S\) 上平滑，且对每个 \(s\in S\)，若 \(\bar s\) 表示 \(\kappa(s)\) 的一个代数闭包的谱，则 \(C_{\bar s}\) 是第 1 节意义下 \(G_{\bar s}\) 的 Cartan 子群。
立即可见，若 \(C\) 是 \(G\) 的 Cartan 子群，则对 \(S\) 上的每个 \(S'\)，\(C_{S'}\) 是 \(G_{S'}\) 的 Cartan 子群。若 \(S\) 是代数闭域的谱，这恢复了第 1 节所回顾的概念。最后，\(G\) 的群子预概形 \(C\) 为 Cartan 子群这一性质，显然对忠实平坦拟紧拓扑是局部的。
\paragraph{\textbf{定理 3.2。}}
\SGARefTarget{sga3:xii:3.2}
\emph{设 \(S\) 上的 \(G\) 如 \SGARefLink{sga3:xii:3.1}{3.1}，并设 \(G\) 在 \(S\) 上仿射，且具有局部常值约化秩。则映射
\[
  T\longmapsto\operatorname{Centr}_G(T)
\]
诱导从 \(G\) 的极大环面集合到 \(G\) 的 Cartan 子群集合的双射。\footnote{这里给出的证明实际上只对 \(G\) 的闭 Cartan 子群证明了\SGARefLink{sga3:xii:3.2}{定理 3.2}。不过，\SGARefLink{sga3:xii:7.1:item:a}{7.1(a)} 证明了所陈述的\SGARefLink{sga3:xii:3.2}{定理 3.2}，并蕴含 \(G\) 的 Cartan 子群是闭的。}若 \(C\) 对应于 \(T\)，则 \(T\) 是 \(C\) 的唯一极大环面。}
若 \(T\) 是 \(G\) 的极大环面，则由 \SGARefLink{sga3:xi:corollary:5:3}{XI~5.3} 或 \SGARefLink{sga3:xii:6.2}{6.2}，函子 \(\operatorname{Centr}_G(T)\) 可由 \(G\) 的闭平滑子预概形
\[
  C=\operatorname{Centr}_G(T)
\]
表示，由定义可见 \(C\) 是 \(G\) 的 Cartan 子群。此外，\(T\) 显然是 \(C\) 的极大环面，且它是中心的，故由 \SGARefLink{sga3:xii:1.13}{1.13}，它是 \(C\) 的唯一极大环面。因此，从极大环面集合到 Cartan 子群集合的映射 \(T\mapsto\operatorname{Centr}_G(T)\) 是单射；只须再证满射性。
设 \(C\) 是 \(G\) 的 Cartan 子群。我们要证明，对 \(G\) 的某个极大环面 \(T\)，有 \(C=\operatorname{Centr}_G(T)\)。只须找到 \(C\) 的极大环面 \(T\)，因为此时 \(T\) 是 \(G\) 的极大环面：对每个 \(s\in S\)，群 \(G_s\) 与 \(C_s\) 具有相同的约化秩。由 \SGARefLink{sga3:ix:theorem:5:6:item:b}{IX~5.6(b)}，\(T\) 位于 \(C\) 的中心，故
\[
  C\subset C'=\operatorname{Centr}_G(T).
\]
\(S\) 上平滑群 \(C'\) 的平滑子群 \(C\) 与 \(C'\) 逐纤维相同；故 \(C=C'\)。
由于 \(G\) 具有局部常值约化秩，\(C\) 亦然。因此，由 \SGARefLink{sga3:xii:1.7:item:b}{1.7(b)}，\(C\) 对平展拓扑局部地存在极大环面；根据上述论证，这个环面是中心的，于是 \SGARefLink{sga3:xii:1.13}{1.13} 表明 \(C\) 本身存在极大环面。证明完毕。
\paragraph{\textbf{推论 3.3。}}
\SGARefTarget{sga3:xii:3.3}
\emph{设 \(G\) 是 \(S\) 上具有局部常值约化秩的平滑仿射群预概形，令
\[
  \mathcal C:(\mathrm{Sch})^\circ_{/S}\longrightarrow(\mathrm{Ens})
\]
为如下定义的函子
\[
  \mathcal C(S')=
  \{G_{S'}\text{ 的 Cartan 子群}\}.
\]
则 \(\mathcal C\) 同构于 \SGARefLink{sga3:xii:1.10}{1.10} 的函子 \(\mathcal T\)，因而可由同一个在 \(S\) 上平滑、分离且有限型的预概形表示。}
\paragraph{\textbf{推论 3.4。}}
\SGARefTarget{sga3:xii:3.4}
\emph{在 \SGARefLink{sga3:xii:3.2}{3.2} 的假设下，若 \(C=\operatorname{Centr}_G(T)\)，则
\[
  \operatorname{Norm}_G(C)=\operatorname{Norm}_G(T).
\]
}
\paragraph{\textbf{注 3.5。}}
\SGARefTarget{sga3:xii:3.5}
即使 \(G\) 不具有局部常值约化秩，它仍可能存在 Cartan 子群。例如，若 \(G\) 具有连通幂零纤维，则 \(G\) 是自身的 Cartan 子群，但其约化秩不一定局部常值；参见 \SGARefLink{sga3:xii:1.6}{1.6}。在报告 XV 中，我们将使用下一报告的正则元素理论，发展 Cartan 子群理论，既不假设 \(G\) 在 \(S\) 上仿射，也不假设其约化秩局部常值。关于去掉仿射性假设，另见下文第 6、7 节。
